\documentclass[10pt,a4paper]{amsart}
\usepackage[margin=19mm]{geometry}
\usepackage{amsmath,amssymb,mathtools}
\usepackage[scr=rsfs,cal=euler]{mathalfa}
\usepackage{xfrac}
\usepackage[unicode,hidelinks,bookmarks=false]{hyperref}
\newcommand{\ie}{i.e.\ }
\newcommand{\cf}{cf.\ }
\newcommand{\resp}{respectively\ }
\newcommand{\Subsection}{\subsection}
\providecommand{\nobreakdash}{\nobreak\hbox{-}}
\setlength{\emergencystretch}{2em}
\multlinegap0pt
\usepackage{amssymb}
\usepackage{enumerate}
\usepackage{float}
\usepackage{tikz}
\usepackage{xcolor}
\newtheorem{thm}{Theorem}
\providecommand{\diff}{\mathrm{d}}
\title{On the discretization of the object space in inverse problems with application to cryo-electron microscopy}
\author{G. Mordant$^\dagger$, L. Evans$^{*}$, D. Silva-S{\'a}nchez$^{\dagger}$, P. Cossio$^{*,\circ}$, R. Lederman$^{\dagger}$}
\date{Source: arXiv:2609.01688v1 (CC BY 4.0)}
\begin{document}
\maketitle

\noindent\emph{Source and licence.} On the discretization of the object space in inverse problems with application to cryo-electron microscopy. Version arXiv:2609.01688v1 (CC BY 4.0), distributed by arXiv under the Creative Commons Attribution 4.0 International licence, http://creativecommons.org/licenses/by/4.0/, as recorded in the arXiv licence metadata for this exact version and shown on the abstract page https://arxiv.org/abs/2609.01688v1. Full licence notice: \texttt{licenses/arxiv-2609.01688-CC-BY-4.0.txt}.\par\smallskip

\noindent\textit{Editorial scope.} Counted unit: the article's thm LowBd (source0.tex lines 633--646). The article's complete original proof of it is delivered with it (source0.tex lines 1399--1583, one \texttt{proof} environment of 185 lines, which the article places away from the statement). No mathematics is written, repaired or re-derived: the statement and the proof are the article's own lines, in the article's own notation, and no retained line is substituted or re-flowed.  Retained uncounted as the source-local context the proof needs: lines 611--621.  Nothing was omitted from any retained span, so the package carries no omission record.  No retained line contains a citation, so the wrapper carries no bibliography and no bibliography entry.  No retained line contains a figure environment, an image-inclusion command or drawing code, so no image is delivered and the package declares no asset dependency.  Editorial formatting only, in addition: an AMS \texttt{amsart} preamble that restates the article's own declarations and macros used by the retained lines, namely the environments \texttt{thm} on separate counters, together with the standard packages those lines need.  The retained environments print in this document as Theorem 1 in order of appearance, whereas the article prints the same items with the numbering of its own full document.  The complete native source of the article as distributed in the arXiv e-print for this exact version is bundled unchanged as \texttt{sources/209135/source0.tex}.
\begin{equation}
\label{eq: NumMom}
\begin{aligned}
\mathfrak m:=\sup\bigg\{\mathcal K\in\mathbb N_0:\ &
\exists\alpha\in\Delta_{M-1}\ \text{such that}\\[-2pt]
&\sum_{m=1}^M\alpha_m x_m^\kappa
=\int x^\kappa\,\diff\rho(x)
\ \text{for every }|\kappa|\leq\mathcal K
\bigg\}\in\mathbb N_0\cup\{+\infty\}.
\end{aligned}
\end{equation}

\begin{thm}
\label{thm: LowBd}
Suppose that $\rho$ has bounded support and let $\mathfrak m<\infty$ be
defined by~\eqref{eq: NumMom}. Then, for constants $c>0$ and
$\sigma_0<\infty$ depending only on $\rho$ and the candidate points,
\begin{equation}
\label{eq:uniform-kl-lower-bound}
\inf_{\alpha\in\Delta_{M-1}}
\operatorname{KL}(p_\sigma*\rho\|p_\sigma*\mathrm P_M^\alpha)
\geq
c\,\sigma^{-2(\mathfrak m+1)},
\qquad \sigma\geq\sigma_0.
\end{equation}
\end{thm}

\begin{proof}[Proof of Theorem~\ref{thm: LowBd}]

For \(1\leq k\leq \mathfrak{m}+1\), define the symmetric
moment-difference tensor
\[
  T_k(\alpha)
  :=
  \int x^{\otimes k}\,
  d\bigl(\mathrm{P}_M^\alpha-\rho\bigr)(x).
\]
Matching all coordinate moments of total degree \(k\) is equivalent to
\(T_k(\alpha)=0\). By maximality of \(\mathfrak{m}\), no
\(\alpha\in\Delta_{M-1}\) makes
\[
  T_1(\alpha),\ldots,T_{\mathfrak{m}+1}(\alpha)
\]
all zero. These tensors depend continuously on \(\alpha\), and the
simplex is compact. Consequently,
\begin{equation}
  \delta
  :=
  \min_{\alpha\in\Delta_{M-1}}
  \left\{
    \sum_{k=1}^{\mathfrak{m}+1}
    \|T_k(\alpha)\|_F^2
  \right\}^{1/2}
  >0.
  \label{eq:uniform-moment-gap}
\end{equation}

Let
\[
  \Delta_\alpha(u)
  :=
  \mathcal{F}\bigl[\mathrm{P}_M^\alpha\bigr](u)
  -
  \mathcal{F}[\rho](u),
\]
where, for a finite measure \(\mu\),
\[
  \mathcal{F}[\mu](u)
  :=
  \int e^{i\langle u,x\rangle}\,d\mu(x).
\]
The supports of \(\rho\) and of every \(\mathrm{P}_M^\alpha\) lie in
one fixed ball. Taylor's formula for \(e^{i\langle u,x\rangle}\)
therefore gives, uniformly in \(\alpha\) and for \(\|u\|\leq 1\),
\begin{equation}
  \Delta_\alpha(u)
  =
  \sum_{k=1}^{\mathfrak{m}+1}
  \frac{i^k}{k!}
  \left\langle T_k(\alpha),u^{\otimes k}\right\rangle
  +
  E_\alpha(u),
  \qquad
  |E_\alpha(u)|
  \leq
  C_1\|u\|^{\mathfrak{m}+2}.
  \label{eq:uniform-characteristic-expansion}
\end{equation}

Let $B_1:=\{z\in\mathbb R^d:\|z\|\leq1\}$. The map
\[
  (A_1,\ldots,A_{\mathfrak{m}+1})
  \longmapsto
  \sum_{k=1}^{\mathfrak{m}+1}
  \frac{i^k}{k!}
  \left\langle A_k,z^{\otimes k}\right\rangle
\]
is injective on
\(
  \bigoplus_{k=1}^{\mathfrak{m}+1}
  \operatorname{Sym}^k(\mathbb{R}^d).
\)
Hence the \(L^2(B_1)\)-norm of the resulting polynomial is equivalent
to the Euclidean norm of its coefficients. It follows from
\eqref{eq:uniform-moment-gap} that, for every \(\sigma\geq 1\),
\begin{align}
  &\left\|
    \sum_{k=1}^{\mathfrak{m}+1}
    \frac{i^k\sigma^{-k}}{k!}
    \left\langle T_k(\alpha),z^{\otimes k}\right\rangle
  \right\|_{L^2(B_1)}
\geq
  c_1
  \left\{
    \sum_{k=1}^{\mathfrak{m}+1}
    \sigma^{-2k}\|T_k(\alpha)\|_F^2
  \right\}^{1/2}
  \geq
  c_1\delta\,
  \sigma^{-(\mathfrak{m}+1)}.
  \label{eq:polynomial-norm-lower-bound}
\end{align}

Set \(u=z/\sigma\). By
\eqref{eq:uniform-characteristic-expansion}, the \(L^2(B_1)\)-norm of
the remainder satisfies
\[
  \bigl\|E_\alpha(z/\sigma)\bigr\|_{L^2(B_1)}
  \leq
  C_2\sigma^{-(\mathfrak{m}+2)}
\]
uniformly in \(\alpha\). Therefore,
\begin{align*}
  \bigl\|\Delta_\alpha(z/\sigma)\bigr\|_{L^2(B_1)}
  &\geq
  c_1\delta\,\sigma^{-(\mathfrak{m}+1)}
  -
  C_2\sigma^{-(\mathfrak{m}+2)}.
\end{align*}
Choosing
\(
  \sigma_0
  \geq
  \max\left\{
    1,\frac{2C_2}{c_1\delta}
  \right\}\), it holds , for every \(\sigma\geq\sigma_0\), uniformly in
\(\alpha\in\Delta_{M-1}\), that
\[
  \bigl\|\Delta_\alpha(z/\sigma)\bigr\|_{L^2(B_1)}
  \geq
  \frac{c_1\delta}{2}
  \sigma^{-(\mathfrak{m}+1)}.
\]
After squaring, this gives
\begin{equation}
  \int_{B_1}
  \left|\Delta_\alpha(z/\sigma)\right|^2\,dz
  \geq
  c_2\sigma^{-2(\mathfrak{m}+1)}.
  \label{eq:characteristic-lower-bound}
\end{equation}

Since
\[
  \mathcal{F}[p_\sigma](u)
  =
  e^{-\sigma^2\|u\|^2/2},
\]
Plancherel's identity and the change of variables \(z=\sigma u\)
yield
\begin{align}
  \left\|
    p_\sigma*\mathrm{P}_M^\alpha-p_\sigma*\rho
  \right\|_2^2
  &=
  c_d
  \int_{\mathbb{R}^d}
  e^{-\sigma^2\|u\|^2}
  |\Delta_\alpha(u)|^2\,du
  \notag\\
  &=
  c_d\sigma^{-d}
  \int_{\mathbb{R}^d}
  e^{-\|z\|^2}
  |\Delta_\alpha(z/\sigma)|^2\,dz
  \notag\\
  &\geq
  c_de^{-1}\sigma^{-d}
  \int_{B_1}
  |\Delta_\alpha(z/\sigma)|^2\,dz
  \notag\\
  &\geq
  c_3
  \sigma^{-d-2(\mathfrak{m}+1)}.
  \label{eq:smoothed-l2-lower-bound}
\end{align}


Both $p_\sigma*\rho$ and $p_\sigma*\mathrm P_M^\alpha$ are bounded above by
$B_\sigma:=(2\pi\sigma^2)^{-d/2}$. The KL--Hellinger inequality gives
\[
\operatorname{KL}(p_\sigma*\rho\|p_\sigma*\mathrm P_M^\alpha)
\geq\frac1{4B_\sigma}
\|(p_\sigma*\rho)-(p_\sigma*\mathrm P_M^\alpha)\|_2^2.
\]
Combining this with~\eqref{eq:smoothed-l2-lower-bound} gives, uniformly in
$\alpha$,
\(
\operatorname{KL}(p_\sigma*\rho\|p_\sigma*\mathrm P_M^\alpha)
\geq c\,\sigma^{-2(\mathfrak m+1)}.
\)
\end{proof}

\end{document}
